\part{Training for context}

\chapter{What training would add}
\label{ch:training-adds}
Training is not the project goal, but knowing what it would add clarifies the inference focus.



\begin{intuitionbox}{Intuition}
Training turns forward execution into differentiable optimisation.
\end{intuitionbox}

\begin{definitionbox}{Mathematical or contract statement}
Training requires gradients, loss, optimiser state, dataset pipelines, checkpoint writing, and often distributed execution.
\end{definitionbox}

\begin{implementationbox}{Implementation interpretation}
Add backward methods, gradient buffers, optimiser hooks, and training checkpoint formats only in a separate project phase.
\end{implementationbox}

\begin{verificationbox}{How to verify this works}
\begin{itemize}
\item Check shapes and dtypes at the boundary.
\item Compare deterministic outputs against PyTorch golden vectors.
\item Record tolerance, seed, model artifact, and backend.
\item Do not benchmark until validation passes.
\end{itemize}
\end{verificationbox}

\begin{warningbox}{Common failure modes}
\begin{itemize}
\item Letting framework defaults choose dtype or layout.
\item Testing only a batch size of one.
\item Depending on upstream checkpoint names in runtime code.
\item Treating benchmark output as validation.
\end{itemize}
\end{warningbox}

\begin{chapterexercises}
\item State the central invariant in one sentence.\\
\textit{Hint.} Use the conceptual guide vocabulary.
\item Construct a toy example with small shapes.\\
\textit{Hint.} Use $B=2$, $T=3$, $H=8$ where possible.
\item Name a validation test that would catch a plausible bug.\\
\textit{Hint.} Prefer generated golden vectors to handwritten long arrays.
\end{chapterexercises}

\chapter{Language-modelling loss}
\label{ch:lm-loss}
The standard autoregressive training objective is next-token prediction.

\[
\mathcal{L}(\theta)=-\sum_{t=1}^{T-1}\log p_{\theta}(x_{t+1}\mid x_{\leq t}).
\]

\begin{intuitionbox}{Intuition}
The model is penalised for assigning low probability to the observed next token.
\end{intuitionbox}

\begin{definitionbox}{Mathematical or contract statement}
The cross-entropy loss sums negative log probabilities of next tokens.
\end{definitionbox}

\begin{implementationbox}{Implementation interpretation}
Training implementation would shift logits and labels, mask padding, and compute stable log-softmax.
\end{implementationbox}

\begin{verificationbox}{How to verify this works}
\begin{itemize}
\item Check shapes and dtypes at the boundary.
\item Compare deterministic outputs against PyTorch golden vectors.
\item Record tolerance, seed, model artifact, and backend.
\item Do not benchmark until validation passes.
\end{itemize}
\end{verificationbox}

\begin{warningbox}{Common failure modes}
\begin{itemize}
\item Letting framework defaults choose dtype or layout.
\item Testing only a batch size of one.
\item Depending on upstream checkpoint names in runtime code.
\item Treating benchmark output as validation.
\end{itemize}
\end{warningbox}

\begin{chapterexercises}
\item State the central invariant in one sentence.\\
\textit{Hint.} Use the conceptual guide vocabulary.
\item Construct a toy example with small shapes.\\
\textit{Hint.} Use $B=2$, $T=3$, $H=8$ where possible.
\item Name a validation test that would catch a plausible bug.\\
\textit{Hint.} Prefer generated golden vectors to handwritten long arrays.
\end{chapterexercises}

\chapter{Backpropagation overview}
\label{ch:backprop}
Backpropagation computes derivatives of loss with respect to parameters and activations.



\begin{intuitionbox}{Intuition}
Backward pass reverses the forward graph.
\end{intuitionbox}

\begin{definitionbox}{Mathematical or contract statement}
For \(y=f(g(x))\), gradients use the chain rule.
\end{definitionbox}

\begin{implementationbox}{Implementation interpretation}
A training engine must store or recompute activations and implement backward kernels.
\end{implementationbox}

\begin{verificationbox}{How to verify this works}
\begin{itemize}
\item Check shapes and dtypes at the boundary.
\item Compare deterministic outputs against PyTorch golden vectors.
\item Record tolerance, seed, model artifact, and backend.
\item Do not benchmark until validation passes.
\end{itemize}
\end{verificationbox}

\begin{warningbox}{Common failure modes}
\begin{itemize}
\item Letting framework defaults choose dtype or layout.
\item Testing only a batch size of one.
\item Depending on upstream checkpoint names in runtime code.
\item Treating benchmark output as validation.
\end{itemize}
\end{warningbox}

\begin{chapterexercises}
\item State the central invariant in one sentence.\\
\textit{Hint.} Use the conceptual guide vocabulary.
\item Construct a toy example with small shapes.\\
\textit{Hint.} Use $B=2$, $T=3$, $H=8$ where possible.
\item Name a validation test that would catch a plausible bug.\\
\textit{Hint.} Prefer generated golden vectors to handwritten long arrays.
\end{chapterexercises}

\chapter{Optimisers and state}
\label{ch:optimisers}
Optimisers turn gradients into parameter updates and often keep additional state.



\begin{intuitionbox}{Intuition}
Optimiser state is another model-sized system.
\end{intuitionbox}

\begin{definitionbox}{Mathematical or contract statement}
SGD is \(\theta\leftarrow\theta-\eta\nabla_\theta L\); adaptive methods store moving statistics.
\end{definitionbox}

\begin{implementationbox}{Implementation interpretation}
Training support would require optimiser-state checkpoints separate from inference artifacts.
\end{implementationbox}

\begin{verificationbox}{How to verify this works}
\begin{itemize}
\item Check shapes and dtypes at the boundary.
\item Compare deterministic outputs against PyTorch golden vectors.
\item Record tolerance, seed, model artifact, and backend.
\item Do not benchmark until validation passes.
\end{itemize}
\end{verificationbox}

\begin{warningbox}{Common failure modes}
\begin{itemize}
\item Letting framework defaults choose dtype or layout.
\item Testing only a batch size of one.
\item Depending on upstream checkpoint names in runtime code.
\item Treating benchmark output as validation.
\end{itemize}
\end{warningbox}

\begin{chapterexercises}
\item State the central invariant in one sentence.\\
\textit{Hint.} Use the conceptual guide vocabulary.
\item Construct a toy example with small shapes.\\
\textit{Hint.} Use $B=2$, $T=3$, $H=8$ where possible.
\item Name a validation test that would catch a plausible bug.\\
\textit{Hint.} Prefer generated golden vectors to handwritten long arrays.
\end{chapterexercises}

\chapter{Why inference is a different engineering problem}
\label{ch:inference-different}
Inference focuses on artifact loading, forward computation, cache state, generation, latency, throughput, and reproducibility.



\begin{intuitionbox}{Intuition}
Inference is not simpler; it is differently constrained.
\end{intuitionbox}

\begin{definitionbox}{Mathematical or contract statement}
A good inference engine guarantees semantics, stable loading, cache correctness, and measured performance.
\end{definitionbox}

\begin{implementationbox}{Implementation interpretation}
This project focuses on inference to complete the cross-language contract before training complexity.
\end{implementationbox}

\begin{verificationbox}{How to verify this works}
\begin{itemize}
\item Check shapes and dtypes at the boundary.
\item Compare deterministic outputs against PyTorch golden vectors.
\item Record tolerance, seed, model artifact, and backend.
\item Do not benchmark until validation passes.
\end{itemize}
\end{verificationbox}

\begin{warningbox}{Common failure modes}
\begin{itemize}
\item Letting framework defaults choose dtype or layout.
\item Testing only a batch size of one.
\item Depending on upstream checkpoint names in runtime code.
\item Treating benchmark output as validation.
\end{itemize}
\end{warningbox}

\begin{chapterexercises}
\item State the central invariant in one sentence.\\
\textit{Hint.} Use the conceptual guide vocabulary.
\item Construct a toy example with small shapes.\\
\textit{Hint.} Use $B=2$, $T=3$, $H=8$ where possible.
\item Name a validation test that would catch a plausible bug.\\
\textit{Hint.} Prefer generated golden vectors to handwritten long arrays.
\end{chapterexercises}
